%% kapa: exam=1 round=36 year=2025 period=1 form=A subject=economics subject_ko=경제학원론 seq=1 number=41
%% source: Q-Net 감정평가사 기출문제 (cst00309) artlSeq=5231525 file=raw/1008/5231525/2025년 제36회 감정평가사 1차시험 1교시 문제지 원본.pdf
%% conversion: kordoc 4.24.51 -> tools/gen1.py (auto); review=proofread-sonnet; flags=none
\begin{question}{41}
X재 수요의 탄력성에 관한 설명으로 옳지 않은 것은? (단, 주어진 소득으로 X재와 Y재만 양(+)의 소비를 한다)
\begin{choices}
  \item X재의 가격탄력성이 1일 때, X재 가격의 변동은 X재의 지출액을 변화시키지 않는다.
  \item X재 소득탄력성이 1이면, 소득소비곡선은 원점을 지나는 직선이다.
  \item X재 소득탄력성이 1보다 크면, Y재 소득탄력성은 1보다 작다.
  \item Y재의 X재 가격에 대한 교차탄력성이 1이면, X재의 가격탄력성은 1이다.
  \item X재 소득탄력성이 1이면, Y재 소득탄력성도 1이다.
\end{choices}
\end{question}
\answer{4}
